\section{The native clause intermediate representation}

\subsection{Why clauses are the right granularity}

A proof graph repeatedly uses statements of the form
\[
       \text{when a finite precondition holds, every requested instance of a
       local invariant holds.}
\]
This is precisely the shape of a clause
\([A;y.B(y)]\), read \(A\to\forall y B(y)\).  In the checked Lean source the
structure is
\begin{lstlisting}[language={},caption={The upstream clause datatype, schematically}]
structure Cl (n : Nat) where
  ante : NFm n
  cons : NFm (n + 1)
\end{lstlisting}
The clause calculus stores such clauses in the context.  Its \path{use} rule
proves an instance \(B(t)\) after proving the antecedent \(A\).

A plain assumption \(X\) is represented by \path{Cl.asm X}, whose antecedent is
\(0=0\) and whose consequent is the shifted formula \(X\).  This representation
is important for program compilation: predecessor outputs are inserted as
ordinary assumption clauses, and \(\KPrf.\mathsf{cut}\) can discharge them with
already compiled proofs.

\subsection{Formula discipline}

The raw clause calculus is optimized for \(Q\)-free native formulas: equations
and the native \(D\)-binder.  Arbitrary native formulas are translated by
\path{NFm.unQ}, which rewrites \(Q_y B\) as a \(D\)-formula whose witness
coordinate is unused.  The checked theorems state
\[
\begin{aligned}
\KPrf(\Gamma^{\mathsf{asm}},A)&\iff\NPrf(\Gamma,A)
    &&\text{for \(Q\)-free sequents},\\
\KPrf((\mathsf{unQ}\,\Gamma)^{\mathsf{asm}},\mathsf{unQ}\,A)
    &\iff\NPrf(\Gamma,A)
    &&\text{in general}.
\end{aligned}
\]
Hence a compiler may normalize all exported node formulas to \(Q\)-free form
without loss of provability.

The native language has no primitive conjunction or implication, but this is
not a limitation on expressiveness: the scalar core defines the usual logical
connectives as macros and is bi-interpretable with ordinary iPRA.  At the
compiler level it is often simpler to split a conjunction into several node
outputs or package several equations into one coded Boolean equation.

\subsection{The five rule families}

The calculus has the following operations.

\begin{enumerate}
\item \textbf{Refl}: derive \(t=t\).
\item \textbf{Eval}: the two iterator computation cases, expressed so that an
existing equation about the counter and one about the iterator produce the
next equation.
\item \textbf{Use}: instantiate a clause after proving its antecedent.
\item \textbf{Elim}: open a \(D\)-formula at a fresh witness and continue under
the corresponding clause.
\item \textbf{Intro}: prove a \(D\)-formula from a progress derivation.  Its
start mode is read from the diagonal shape of the matrix.
\end{enumerate}

For framework design, the crucial checked derived operations are weakening,
renaming, substitution, and cut.  These operations are proved by structural
induction on \(\KPrf\) derivations.  Consequently, a proof compiler may treat
them as trusted library combinators rather than re-proving capture-avoidance and
context transport for every node class.

\subsection{Clause-carrying nodes}

\begin{definition}[Native clause node]
Fix a base clause context \(\Gamma\) and a list \(\Delta\) of predecessor
outputs, each a native formula with \(n\) free variables.  A native clause node
consists of:
\begin{enumerate}
\item an output formula \(A\in\NFm(n)\);
\item a proof that \(A\) is \(Q\)-free, or a prescribed \path{unQ}
normalization;
\item an actual local derivation
\[
   h:\KPrf\bigl(\Delta.\mathsf{map}(\Cl.\mathsf{asm})\mathbin{++}\Gamma,A\bigr);
\]
\item optional executable metadata: demand transformers, resource bounds,
trace descriptions, witness-cloud bounds, or observer tolerances.
\end{enumerate}
\end{definition}

The executable metadata does not justify the node.  The field \(h\) does.  The
metadata is used to generate concrete terms and to make the theorem auditable.

\subsection{Normalized node classes}

Most applications can be normalized to six classes.

\paragraph{Equation nodes.}
These export an equation between concrete PR terms.  Examples include tuple
projection, polynomial evaluation, finite table lookup, rational normalization,
substitution, and fixed-dimensional matrix multiplication.  Their proofs use
computation equations, congruence, and algebraic lemmas.

\paragraph{Clause nodes.}
These add reusable implications of the form \([A;y.B(y)]\).  They encode local
invariants, guarded operations, or universal finite families.  The antecedent
is a guard, not an invitation to use an undecidable partial function.

\paragraph{Trace nodes.}
A trace node contains an initializer, PR step, fuel, quantifier-free invariant,
base derivation, and step derivation.  It exports the terminal invariant or a
clause giving every bounded instance.

\paragraph{Search nodes.}
A search node contains a finite bound and a Booleanized predicate.  It exports a
least witness, a complete list, or the equation recording that no witness was
found.  Search correctness is a trace theorem.

\paragraph{Relational nodes.}
These export a finite witness cloud and compatibility edges.  They avoid
constructing a global selector.  Bounded viability may later select a path, or
a stable observer may eliminate the witness.

\paragraph{Finalizer nodes.}
A finalizer converts verified finite evidence into the exact headline formula.
It is the only place where theorem-specific logical packaging should occur.
Keeping finalization separate makes it clear which part of the development is
uniform and which part encodes the mathematical statement.

\subsection{Effects as data, not proof rules}

Each node may carry an effect record
\[
  (\text{precision},\text{fuel},\text{height},\text{denominator},
    \text{gap},\text{cover},\text{exception budget}).
\]
Effects are natural-number terms and finite certificates.  They guide the host
compiler and appear in contracts.  They are not new logical rules.  For
example, a gap effect is a formula proving a positive rational lower bound; a
height effect is a numeral or PR term bounding every candidate code; a cover
effect is a coded finite list together with a bounded covering check.

This design keeps the trusted proof kernel small while permitting a rich
application-level interface.

\subsection{No semantic reflection}

A native clause program is proof-producing in the literal sense: each node
contains a derivation.  The graph compiler never proves
``the checker returned true, therefore the formula is provable.''  It instead
constructs the formula's derivation from the node's local derivation and the
compiled predecessor derivations.  The checker may be used to generate or
validate data, but the proof term remains the authority.

This distinction is especially important for generated arithmetic.  A Python
program may confirm a million finite cases; unless the cases are compiled into
a bounded trace theorem or a closed proof term, that confirmation is not an
iPRA proof.  Conversely, a short \(\KPrf\) derivation may prove all cases even
when the reference program is present only for regression testing.
